\chapter{强化学习游玩雅达利游戏Pong!}
\label{app:pong}

本内容主要参考自 Andrej Karpathy 的博客文章：
\url{https://karpathy.github.io/2016/05/31/rl/}

\section{引言}

4个限制AI发展的因素：

\begin{enumerate}[
  label=（\arabic*）,
  labelindent=\parindent,
  labelwidth=!,
  leftmargin=*,
  labelsep=0.4em,
  parsep=0.5em
]
  \item \textbf{Compute}（算力）：最显而易见的因素，例如摩尔定律、GPU、ASIC；
  \item \textbf{Data}（数据）：以结构化、可直接使用的形式存在的数据，例如 ImageNet；
  \item \textbf{Algorithms}（算法）：研究与创新，例如反向传播（backprop）、CNN、LSTM；
  \item \textbf{Infrastructure}（基础设施）：底层软件栈，例如 Linux、TCP/IP、Git、ROS、PR2、AWS、AMT、TensorFlow 等。
\end{enumerate}

下面使用130行Python，仅用numpy作为依赖，
从零开始用像素和深度神经网络玩一款ATARI游戏（Pong！）。

\section{像素乒乓球（Pong from Pixels）}

\subsection{背景}

\textbf{游戏规则：}游戏中存在两个乒乓球拍和一个乒乓球，
其中一个球拍由玩家控制，另一个由AI控制，双方需要使用球拍接住球，
并让球通过球拍反弹给另一个球员。

\textbf{规则分析：}每一时刻，玩家/AI接收到一个图像帧
（210×160×3字节的数组），然后对球拍做出对应动作：
向上或是向下移动球拍（二元选择）。
每次做出选择后，游戏模拟器会执行相应动作并计算分数：
如果球越过对手方，则获得+1奖励；如果未能击中球，
则获得-1奖励。

\textbf{目标：}通过移动球拍来获取尽可能多的奖励。

\begin{figure}[H]
  \centering
  \begin{tikzpicture}[
    x=1cm,
    y=1cm,
    >=Stealth,
    flow/.style={
      draw=noteBlue,
      very thick,
      -Stealth,
      line cap=round,
      line join=round
    },
    pixframe/.style={
      draw=black!65,
      very thick,
      fill=brown!70!black
    },
    layerlabel/.style={
      font=\bfseries\small,
      text=white
    },
    label/.style={font=\small},
    smalllabel/.style={font=\small}
  ]
    % --- Atari Pong 像素画面，保持与图 1.2 相近的画法 ---
    \begin{scope}[shift={(0,0)}]
      \node[pixframe, minimum width=2.9cm, minimum height=2.9cm] (screen) at (0,0) {};
      \draw[orange!45, line width=1.3pt]
        (-0.73,1.07) rectangle (-0.59,1.34);
      \draw[green!60!black, line width=1.3pt]
        (0.59,1.07) rectangle (0.73,1.34);
      \draw[white, line width=4.5pt]
        ([xshift=0pt,yshift=-0.69cm]screen.north west) --
        ([xshift=0pt,yshift=-0.69cm]screen.north east);
      \fill[orange!20]
        ([xshift=-1.08cm,yshift=-0.95cm]screen.center) rectangle
        ++(0.09,0.37);
      \fill[green!55!black]
        ([xshift=1.00cm,yshift=-1.40cm]screen.center) rectangle
        ++(0.09,0.37);
      \fill[white] ([xshift=-0.15cm,yshift=-0.28cm]screen.center) rectangle ++(0.04,0.04);
      \node[label, anchor=south] at ([yshift=2pt]screen.north) {一帧原始像素};
      \node[smalllabel, anchor=north, align=center] at ([yshift=-2pt]screen.south)
        {白线下方：\\$210$ 像素宽 $\times$ $160$ 像素高};
    \end{scope}

    \draw[flow] (1.85,0) -- (3.05,0);
    \node[label, anchor=south] at (2.45,0) {转为数组};

    % --- 右侧：三层 RGB 像素数组 ---
    \begin{scope}[shift={(3.8,-1.25)}]
      % B layer.
      \fill[blue!14, fill opacity=0.95] (0.92,0.72) rectangle (3.92,3.12);
      \draw[blue!65!black, very thick, rounded corners=2pt] (0.92,0.72) rectangle (3.92,3.12);
      \foreach \x in {1.42,1.92,2.42,2.92,3.42} {
        \draw[blue!35] (\x,0.72) -- (\x,3.12);
      }
      \foreach \y in {1.12,1.52,1.92,2.32,2.72} {
        \draw[blue!35] (0.92,\y) -- (3.92,\y);
      }
      \fill[blue!65!black] (3.55,2.78) rectangle (3.92,3.12);
      \node[layerlabel] at (3.735,2.95) {B};

      % G layer.
      \fill[green!14, fill opacity=0.95] (0.46,0.36) rectangle (3.46,2.76);
      \draw[green!50!black, very thick, rounded corners=2pt] (0.46,0.36) rectangle (3.46,2.76);
      \foreach \x in {0.96,1.46,1.96,2.46,2.96} {
        \draw[green!35!black!35] (\x,0.36) -- (\x,2.76);
      }
      \foreach \y in {0.76,1.16,1.56,1.96,2.36} {
        \draw[green!35!black!35] (0.46,\y) -- (3.46,\y);
      }
      \fill[green!50!black] (3.09,2.42) rectangle (3.46,2.76);
      \node[layerlabel] at (3.275,2.59) {G};

      % R layer.
      \fill[red!12] (0,0) rectangle (3.0,2.4);
      \draw[red!70!black, very thick, rounded corners=2pt] (0,0) rectangle (3.0,2.4);
      \foreach \x in {0.5,1.0,1.5,2.0,2.5} {
        \draw[red!35] (\x,0) -- (\x,2.4);
      }
      \foreach \y in {0.4,0.8,1.2,1.6,2.0} {
        \draw[red!35] (0,\y) -- (3.0,\y);
      }
      \fill[red!70!black] (2.63,2.06) rectangle (3.0,2.4);
      \node[layerlabel] at (2.815,2.23) {R};

      \node[font=\bfseries\small] at (2.18,3.72) {RGB 三层像素数组};
      \node[label, align=center] at (1.5,-0.55)
        {$210 \times 160 \times 3~\mathrm{Byte}$\\宽度 $\times$ 高度 $\times$ 颜色通道};

      \draw[flow] (4.20,1.82) -- (4.70,1.82);
      \node[smalllabel, align=left, anchor=west] at (4.85,1.82)
        {每个通道元素：\\$1~\mathrm{Byte}=8~\mathrm{bit}$\\取值 $0 \sim 255$（256 种）};
    \end{scope}
  \end{tikzpicture}
  \caption{Pong 图像帧可以表示为一个 $210 \times 160 \times 3~\mathrm{Byte}$ 的 RGB 像素数组。}
  \label{fig:pong-frame-array}
\end{figure}


\subsection{实现}

\subsubsection{策略网络（智能体）}

智能体部分用\textbf{策略网络}（Policy Network）来实现。

策略网络由一个两层神经网络构成，该网络的输入为游戏当前像素（210×160×3字节的数组），
输出采取的动作（向上或向下移动）。

\begin{figure}[H]
  \centering
  \begin{tikzpicture}[
    x=1cm,
    y=1cm,
    >=Stealth,
    font=\small,
    netline/.style={draw=black!75, line width=0.6pt},
    pixframe/.style={
      draw=black!65,
      very thick,
      fill=brown!70!black
    },
    hidden/.style={
      circle,
      draw=black!65,
      line width=0.8pt,
      fill=noteBlue!18,
      minimum size=0.62cm
    },
    output/.style={
      circle,
      draw=black!65,
      line width=0.8pt,
      fill=red!22,
      minimum size=0.62cm
    },
    label/.style={font=\small}
  ]
    % --- Raw pixels ---
    \begin{scope}[shift={(1.775,1.775)}]
      \node[pixframe, minimum width=3.55cm, minimum height=3.55cm] (screen) at (0,0) {};
      \draw[orange!45, line width=1.3pt]
        (-0.89,1.31) rectangle (-0.72,1.64);
      \draw[green!60!black, line width=1.3pt]
        (0.72,1.31) rectangle (0.89,1.64);
      \draw[white, line width=5pt]
        ([xshift=0pt,yshift=-0.845cm]screen.north west) --
        ([xshift=0pt,yshift=-0.845cm]screen.north east);
      \fill[orange!20]
        ([xshift=-1.32cm,yshift=-1.16cm]screen.center) rectangle
        ++(0.11,0.45);
      \fill[green!55!black]
        ([xshift=1.22cm,yshift=-1.71cm]screen.center) rectangle
        ++(0.11,0.45);
      \fill[white] ([xshift=-0.18cm,yshift=-0.34cm]screen.center) rectangle ++(0.05,0.05);
      \node[label, anchor=south] at ([yshift=2pt]screen.north) {原始像素};
    \end{scope}

    % --- Hidden layer nodes ---
    \node[label, text=noteBlue] at (5.65,3.78) {隐藏层};
    \foreach \i/\y in {1/3.25,2/2.50,3/1.75,4/1.00,5/0.25} {
      \node[hidden] (h\i) at (5.55,\y) {};
    }

    \foreach \x/\y in {
      3.55/3.55,
      3.55/2.75,
      3.55/2.05,
      3.55/1.35,
      3.55/0.70,
      3.55/0.05
    } {
      \foreach \i in {1,2,3,4,5} {
        \draw[netline] (\x,\y) -- (h\i.west);
      }
    }

    % Output node and connections.
    \node[output] (out) at (7.85,1.75) {};
    \foreach \i in {1,2,3,4,5} {
      \draw[netline] (h\i.east) -- (out.west);
    }

    \node[label, text=red!70!black, align=left, anchor=west]
      at (8.18,2.35) {向上或向下\\的概率};
  \end{tikzpicture}
  \caption{两层全连接策略网络从原始像素输出向上或向下的概率}
  \label{fig:pong-two-layer-policy-network}
\end{figure}


代码部分如下：

\begin{lstlisting}[language=Python, aboveskip=\baselineskip, belowskip=\baselineskip]
h = np.dot(W1, x) # 计算隐藏层神经元激活值
h[h < 0] = 0 # ReLU非线性：将小于0的值截断为0
logp = np.dot(W2, h) # 计算向上移动的对数概率
p = 1.0 / (1.0 + np.exp(-logp)) # sigmoid函数：得到向上移动的概率
\end{lstlisting}

$x$作为输入，包含了游戏画面的像素信息；$w_1$和$w_2$分别为连接输入向量$x$和隐藏层$h$、隐藏层$h$
和动作概率$p$各节点的权重，没有使用偏置$b$；最终的输出$p$固定为向上移动球拍的概率（$1 - p$就是
向下移动球拍的概率）。

因为隐藏层没有输出数值大小限制，使用ReLU函数；而最终的输出是一个概率，要求处于$[0,1]$之间，使用
sigmoid函数。

\begin{figure}[H]
  \centering
  \begin{tikzpicture}[
    x=1cm,
    y=1cm,
    >=Stealth,
    font=\small,
    flow/.style={
      draw=noteBlue,
      very thick,
      -Stealth,
      line cap=round,
      line join=round
    },
    netline/.style={
      draw=black!55,
      line width=0.32pt,
      opacity=0.55
    },
    pixframe/.style={
      draw=black!65,
      very thick,
      fill=brown!70!black
    },
    hidden/.style={
      circle,
      draw=black!65,
      line width=0.75pt,
      fill=noteBlue!16,
      minimum size=0.48cm
    },
    output/.style={
      circle,
      draw=black!65,
      line width=0.75pt,
      fill=red!20,
      minimum size=0.58cm
    },
    label/.style={font=\small},
    graphaxis/.style={draw=black!45, line width=0.35pt, -Stealth},
    channelcell/.style={line width=0.38pt},
    formula/.style={font=\small, fill=white, inner sep=1.5pt}
  ]
    % --- Raw Pong frame. ---
    \begin{scope}[shift={(0,0)}]
      \node[pixframe, minimum width=2.25cm, minimum height=2.25cm] (screen) at (0,0) {};
      \draw[orange!45, line width=1.05pt]
        (-0.56,0.83) rectangle (-0.46,1.04);
      \draw[green!60!black, line width=1.05pt]
        (0.46,0.83) rectangle (0.56,1.04);
      \draw[white, line width=3.4pt]
        ([xshift=0pt,yshift=-0.535cm]screen.north west) --
        ([xshift=0pt,yshift=-0.535cm]screen.north east);
      \fill[orange!20]
        ([xshift=-0.84cm,yshift=-0.74cm]screen.center) rectangle
        ++(0.07,0.29);
      \fill[green!55!black]
        ([xshift=0.78cm,yshift=-1.09cm]screen.center) rectangle
        ++(0.07,0.29);
      \fill[white] ([xshift=-0.12cm,yshift=-0.23cm]screen.center) rectangle ++(0.034,0.034);
      \node[label, anchor=south] at ([yshift=2pt]screen.north) {原始 Pong 帧};
    \end{scope}

    \draw[flow] (1.42,0) -- (2.55,0);
    \node[label, anchor=south] at (1.98,0.20) {预处理};
    \node[label, anchor=north] at (1.98,-0.20) {展平};

    % --- Input vector x: one vertical vector split into RGB blocks. ---
    \begin{scope}[shift={(3.35,-2.25)}]
      \node[label, anchor=south] at (0.69,4.68) {输入向量 $x$};

      \fill[red!10] (0,3.00) rectangle (1.38,4.50);
      \draw[red!65!black, channelcell] (0,3.00) rectangle (1.38,4.50);
      \foreach \y in {3.50,4.00} {
        \draw[red!35] (0,\y) -- (1.38,\y);
      }
      \node[text=red!65!black, anchor=east] at (-0.16,3.75) {R};
      \node at (0.69,4.25) {$x_{R,1}$};
      \node at (0.69,3.75) {$x_{R,2}$};
      \node at (0.69,3.25) {$\vdots$};

      \fill[green!10] (0,1.50) rectangle (1.38,3.00);
      \draw[green!50!black, channelcell] (0,1.50) rectangle (1.38,3.00);
      \foreach \y in {2.00,2.50} {
        \draw[green!35!black!35] (0,\y) -- (1.38,\y);
      }
      \node[text=green!45!black, anchor=east] at (-0.16,2.25) {G};
      \node at (0.69,2.75) {$x_{G,1}$};
      \node at (0.69,2.25) {$x_{G,2}$};
      \node at (0.69,1.75) {$\vdots$};

      \fill[blue!10] (0,0) rectangle (1.38,1.50);
      \draw[blue!65!black, channelcell] (0,0) rectangle (1.38,1.50);
      \foreach \y in {0.50,1.00} {
        \draw[blue!35] (0,\y) -- (1.38,\y);
      }
      \node[text=blue!65!black, anchor=east] at (-0.16,0.75) {B};
      \node at (0.69,1.25) {$x_{B,1}$};
      \node at (0.69,0.75) {$x_{B,2}$};
      \node at (0.69,0.25) {$\vdots$};
    \end{scope}

    % --- Hidden layer. ---
    \node[label, text=noteBlue] at (7.45,1.60) {隐藏层 $h$};
    \foreach \i/\y in {1/1.03,2/0.52,3/0.01,4/-0.50,5/-1.01} {
      \node[hidden] (h\i) at (7.45,\y) {};
    }

    \foreach \y in {1.90,0.95,0.00,-0.95,-1.90} {
      \foreach \i in {1,2,3,4,5} {
        \draw[netline] (4.73,\y) -- (h\i.west);
      }
    }

    \node[formula] at (7.45,-2) {$h=\mathrm{ReLU}(W_1x)$};
    \node[formula] at (5.92,1.88) {$W_1$};

    % --- ReLU function graph, separated from the main connections. ---
    \begin{scope}[shift={(7.45,2.50)}, x=0.55cm, y=0.55cm]
      \draw[graphaxis] (-1.06,0) -- (1.18,0);
      \draw[graphaxis] (0,-0.12) -- (0,1.12);
      \draw[noteOrange, very thick] (-1.0,0) -- (0,0) -- (1.0,1.0);
      \node[label, anchor=south] at (0.32,1.04) {$\mathrm{ReLU}(z)$};
    \end{scope}

    % --- Output probability. ---
    \node[output] (p) at (10.85,0.01) {$p$};
    \foreach \i in {1,2,3,4,5} {
      \draw[netline] (h\i.east) -- (p.west);
    }
    \node[formula] at (10.85,-1) {$p=\sigma(W_2h)$};
    \node[formula] at (9.25,1.25) {$W_2$};
    \node[label, align=left, anchor=west] at (11.25,0.01) {动作概率};

    % --- Sigmoid function graph, separated from the main connections. ---
    \begin{scope}[shift={(10.85,2.50)}, x=0.55cm, y=0.55cm]
      \draw[graphaxis] (-2.10,0) -- (2.20,0);
      \draw[graphaxis] (0,-0.10) -- (0,1.14);
      \draw[black!25, line width=0.3pt] (-2.00,1.00) -- (2.05,1.00);
      \draw[red!70!black, very thick, domain=-2:2, samples=80, smooth]
        plot (\x,{1/(1+exp(-3*\x))});
      \node[label, anchor=south] at (0.62,1.05) {$\sigma(z)$};
    \end{scope}
  \end{tikzpicture}
  \caption{策略网络前向计算示意。}
  \label{fig:pong-preprocessed-pixel-vector}
\end{figure}


从设计上来看，权重$w_1$使隐藏层$h$能够识别游戏当前形势
（球在哪、拍在哪），而权重$W_2$影响输出$p$，决定球拍向上移动还是向下移动。

随机权重矩阵$w_1$、$w_2$的初始值，智能体控制球拍会随机乱动。
为了让智能体学会打乒乓球，我们需要做的就是调整权重。

\subsubsection{预处理}

目前我们设计的网络输入是当前的一帧图片，如果只输入当前帧到网络，
决策没有考虑到球拍和球的运动。所以希望向策略网络输入至少两帧图像，
这样它才能检测到运动。

然而这样存在\textbf{信用分配问题}，奖励并不是实时反馈的，
你甚至无法预知反馈会在多少个时间步后返回。当前时刻对球拍的
移动无法知道在多久后导致游戏的胜利；换一个角度说，如果当前时刻
取得了胜利，很难分辨出是之前哪一步的动作导致了胜利，更分不清之前的
哪些动作的贡献大小是怎么分布的。

\subsubsection{监督学习}

我们可以参考监督学习的方法解决问题。

在监督学习中，我们将图像输入网络，输出两个类别UP（上移）和DOWN（下移）的概率。
假如UP的对数概率为$-1.2$，DOWN的对数概率为$-0.36$。（概率分别为$30\%$ 和 $70\%$）

在监督学习中我们可以拿到标签。例如，我们可能会被告知，
当前正确的动作是 UP（标签 0）。
在实现中，我们会将 $+1.0$ 的梯度填入 UP 的对数概率位置，
然后运行反向传播，计算梯度向量
$\nabla_{\params}\log p(y=\mathrm{UP}\mid x)$。

这个梯度会告诉我们：应当如何调整网络中的每一个参数，
才能使网络在遇到类似图像时略微更倾向于预测 UP。
例如，假设网络中有一百万个参数中的某一个梯度为 $-2.1$，
它的含义是：若将该参数增大一个小的正值（例如 $0.001$），
那么 UP 的对数概率就会下降 $2.1 \times 0.001$
（下降是因为符号为负）。
随后我们进行一次参数更新，棒！我们的网络在将来再看到非常相似的图像时，
就会略微更倾向于预测 UP。

\begin{figure}[H]
  \centering
  \begin{tikzpicture}[
    x=1cm,
    y=1cm,
    >=Stealth,
    font=\small,
    flow/.style={
      draw=noteBlue,
      line width=1.05pt,
      -Stealth,
      line cap=round,
      line join=round
    },
    backflow/.style={
      draw=red!75!black,
      line width=1.05pt,
      dashed,
      -Stealth,
      line cap=round,
      line join=round
    },
    netline/.style={
      draw=black!55,
      line width=0.75pt,
      -Stealth,
      line cap=round,
      line join=round
    },
    pixframe/.style={
      draw=black!65,
      very thick,
      fill=brown!70!black
    },
    compute/.style={
      draw=noteBlue!75!black,
      line width=1.0pt,
      rounded corners=4pt,
      fill=noteBlue!6,
      minimum width=2.95cm,
      minimum height=1.30cm,
      align=center,
      inner sep=6pt
    },
    output/.style={
      draw=black!65,
      line width=0.85pt,
      rounded corners=3pt,
      fill=red!10,
      minimum width=1.55cm,
      minimum height=0.58cm,
      align=center,
      inner sep=3pt
    },
    target/.style={
      draw=black!55,
      line width=0.75pt,
      rounded corners=2pt,
      fill=white,
      minimum width=0.58cm,
      minimum height=0.50cm,
      align=center,
      inner sep=2pt
    },
    label/.style={font=\small}
  ]
    \begin{scope}[shift={(0,0)}]
      \node[pixframe, minimum width=2.25cm, minimum height=2.25cm] (screen) at (0,0) {};
      \draw[orange!45, line width=1.05pt]
        (-0.55,0.83) rectangle (-0.44,1.05);
      \draw[green!60!black, line width=1.05pt]
        (0.44,0.83) rectangle (0.55,1.05);
      \draw[white, line width=3.4pt]
        ([yshift=-0.54cm]screen.north west) --
        ([yshift=-0.54cm]screen.north east);
      \fill[orange!20]
        ([xshift=-0.84cm,yshift=-0.73cm]screen.center) rectangle ++(0.07,0.30);
      \fill[green!55!black]
        ([xshift=0.78cm,yshift=-1.05cm]screen.center) rectangle ++(0.07,0.30);
      \fill[white]
        ([xshift=-0.09cm,yshift=-0.20cm]screen.center) rectangle ++(0.035,0.035);
      \node[label, anchor=south] at ([yshift=3pt]screen.north) {输入图像 $x$};
    \end{scope}

    \node[compute] (block) at (3.55,0)
      {可微分计算块\\{\scriptsize 例如神经网络}};

    \draw[flow]
      ([xshift=0.05cm,yshift=0.36cm]block.north west) --
      ([xshift=-0.05cm,yshift=0.36cm]block.north east);
    \node[label, text=noteBlue, anchor=south]
      at ([yshift=0.41cm]block.north) {前向传播};

    \draw[backflow]
      ([xshift=-0.05cm,yshift=-0.36cm]block.south east) --
      ([xshift=0.05cm,yshift=-0.36cm]block.south west);
    \node[label, text=red!75!black, anchor=north]
      at ([yshift=-0.41cm]block.south) {反向传播};

    \draw[netline] (screen.east) -- node[label, above, midway] {$x$} (block.west);

    \node[output, draw=noteGreen!55!black, fill=noteGreen!16] (up) at (6.20,0.58) {向上移动};
    \node[output] (down) at (6.20,-0.58) {向下移动};
    \draw[netline] (block.east) -- (up.west);
    \draw[netline] (block.east) -- (down.west);

    \node[target, draw=noteGreen!55!black] (up-target) at (7.80,0.58) {$1$};
    \node[target] (down-target) at (7.80,-0.58) {$0$};
    \node[label, anchor=south] at (7.80,1.02) {标签};

    \draw[netline, shorten >=2pt] (up.east) -- (up-target.west);
    \draw[netline, shorten >=2pt] (down.east) -- (down-target.west);
  \end{tikzpicture}
  \caption{输入图像经过可微分计算块得到两个动作候选；监督学习标签将正确的向上移动记为 $1$，向下移动记为 $0$。}
  \label{fig:pong-supervised-learning}
\end{figure}




\section{代码分析}
